\chapter{Time Systems and Conversions}
\label{chap:time-conversions}

This chapter fixes the mathematical contract for \lupnt{}'s time systems: the
epoch origin, the scale definitions, the exact offset/rate/periodic equations,
the body and reference system each scale is proper to, and the leap-second and
Earth-orientation inputs. Every conversion is tied to the function that
implements it in \file{cpp/lupnt/conversions/time_conversions.h} and
\file{cpp/lupnt/conversions/epoch.h}; the Python surface is
\py{pnt.convert\_time} and friends from
\file{python/bindings/py_time_converter.cc}.

The relativistic content follows \citep[Ch.~2.3]{iiyama2026thesis} and the
lunar time-scale formulation of \citet{turyshev2026cislunar}, with the
barycentric relation taken from the DE440 ephemeris paper
\citep{park2021de440}. This chapter states what the implementation is
supposed to compute.

\section{The unit and argument contract}
\label{sec:tc-contract}

Every scalar time in the API is a \emph{coordinate offset in seconds from the
J2000 epoch of the corresponding scale}, not a calendar date. The J2000 origin
is
%
\begin{equation}
\label{eq:tc-j2000}
  \mathrm{MJD}_{\mathrm{J2000}} = 51544.5\ \text{(TT days)},
  \qquad
  \mathrm{JD}_{\mathrm{J2000}}  = 2451545.0\ \text{(TT days)} .
\end{equation}
%
The Julian and Modified Julian bridges are exact linear maps,
%
\begin{equation}
\label{eq:tc-mjd}
  t = \bigl(\mathrm{MJD} - 51544.5\bigr)\cdot 86400,
  \qquad
  \mathrm{MJD} = t/86400 + 51544.5 ,
\end{equation}
%
implemented as one-liners:

\begin{lstlisting}[language=cpp,caption={Julian-date bridges.},label={lst:tc-mjd}]
// cpp/lupnt/conversions/time_conversions.cc :: MjdToTime / TimeToMjd
Real MjdToTime(Real mjd) { return (mjd - MJD_J2000_TT) * SECS_DAY; }
Real TimeToMjd(Real t)   { return t / SECS_DAY + MJD_J2000_TT; }
\end{lstlisting}

The \code{Time} enum in \file{cpp/lupnt/core/constants.h} enumerates the ten
supported scales:
%
\begin{equation*}
  \code{UT1},\ \code{UTC},\ \code{TAI},\ \code{TDB},\ \code{TT},\
  \code{TCG},\ \code{TCB},\ \code{GPS},\ \code{TCL},\ \code{LT}.
\end{equation*}

\begin{lupntnote}
A Julian Date is a \emph{representation} of an instant rather than a time
scale, so it is not a member of the enum. Use \cpp{JdToTime} / \cpp{TimeToJd}
together with the relevant scale.
\end{lupntnote}

\section{The precision problem and the \texttt{Epoch} class}
\label{sec:tc-epoch}

\subsection{Why a bare \texttt{double} is not enough}

The natural way to store an instant is a single floating-point number of
seconds elapsed since J2000. A binary64 value \citep{ieee754} carries a fixed
52-bit fraction, so its resolution scales with its own magnitude: the gap
between neighbouring representable values --- one unit in the last place (ULP)
--- is
%
\begin{equation}
\label{eq:tc-ulp}
  \mathrm{ULP}(t) = |t|\,2^{-52} .
\end{equation}
%
At present-day dates $|t| \sim \SI{e9}{\second}$, and \cref{eq:tc-ulp} gives
%
\begin{equation}
\label{eq:tc-ulp-value}
  \mathrm{ULP} \approx \SI{245}{\nano\second} \approx \SI{74}{\meter}/c .
\end{equation}
%
No instant can be written down more precisely than that, and neither can any
quantity derived from one: a time-scale offset, a light time, or an
integration step is snapped to the same grid. For a navigation library,
\SI{74}{\meter} of equivalent range is a hard floor that no improvement to the
physical models can lift.

\subsection{The \texttt{Epoch} representation}

\cpp{Epoch} removes the floor by never storing the large and the small parts
of an instant in the same floating-point number. An epoch is a triple
%
\begin{equation}
\label{eq:tc-epochrep}
  \code{Epoch} =
  \bigl(\, \underbrace{s}_{\text{\code{int64}, whole seconds from J2000}},\;\;
           \underbrace{f}_{\code{Real},\; f \in [0,1)},\;\;
           \underbrace{S}_{\text{time scale}} \,\bigr),
  \qquad t = s + f .
\end{equation}
%
The whole seconds live in a 64-bit integer, which is exact: an \code{int64}
counts every second for $\pm 2.9\times10^{11}$ years without rounding. Only the
sub-second remainder $f$ is floating-point, and because $|f| < 1$ its ULP is
$2^{-52} \approx \SI{0.2}{\femto\second}$, independent of the date. The time
scale $S$ is carried with the value, so subtracting a TT epoch from a TDB one
is rejected rather than silently returning a wrong answer.

Two operations then keep full precision:
%
\begin{align}
  \label{eq:tc-epochdiff}
  \text{difference:}\quad & t_1 - t_2 = \underbrace{(s_1 - s_2)}_{\text{exact, integer}}
                              + (f_1 - f_2), \\
  \label{eq:tc-epochconv}
  \text{conversion:}\quad & \code{Epoch}(s,\, f + \Delta_{A\to B},\, B).
\end{align}
%
In \cref{eq:tc-epochdiff} the integer parts cancel exactly, so the result is
carried at the precision of the \emph{difference} rather than of the epochs:
two epochs a second apart differ to ${\sim}\SI{e-16}{\second}$. In
\cref{eq:tc-epochconv} the inter-scale offset $\Delta_{A\to B}$ never exceeds
${\sim}\SI{100}{\second}$ across the supported scales, so adding it to $f$ is
good to ${\sim}\SI{e-14}{\second}$. Neither operation ever subtracts two large
numbers.

\begin{lupntnote}
\textbf{Autodiff contract.} Derivatives live entirely on $f$; $s$ is a plain
integer and is never differentiated. This matches how time enters estimation:
the epoch is a fixed parameter, while the estimated time-like quantities
(clock bias, clock drift) are exactly the ones that stay in $f$.
Renormalisation shifts whole seconds between $s$ and $f$, which is
derivative-preserving because it subtracts an integer constant. See
\cref{chap:autodiff}.
\end{lupntnote}

\begin{lstlisting}[language=cpp,caption={The split representation and its exact difference.},label={lst:tc-epoch}]
// cpp/lupnt/conversions/epoch.h :: Epoch
class Epoch {
 public:
  Epoch(int64_t sec, Real frac, Time scale) : sec_(sec), frac_(frac), scale_(scale) {
    Normalize();          // fold whole seconds out of frac_, so frac_ in [0,1)
  }
  static Epoch FromSeconds(Real t, Time scale);
  static Epoch FromGregorian(int year, int month, int day, int hour, int min,
                             Real sec, Time scale);

  Real ToSeconds() const {              // LOSSY: reintroduces the 245 ns ULP
    return Real(static_cast<double>(sec_)) + frac_;
  }
  Epoch To(Time target) const;          // offset-composing conversion
  Real operator-(const Epoch& other) const;   // exact: integer parts cancel

 private:
  int64_t sec_ = 0;     // exact
  Real    frac_ = 0.0;  // the only differentiated component
  Time    scale_ = Time::TAI;
};
\end{lstlisting}

\subsection{Collapsing to a \texttt{Real}, deliberately}
\label{sec:tc-collapse}

\cpp{Epoch::ToSeconds()} returns $s + f$ as a single \code{Real}. This is lossy
by construction --- it re-imposes the \SI{245}{\nano\second} quantum of
\cref{eq:tc-ulp} --- and it exists because most of the library's interfaces
predate \cpp{Epoch} and take a bare seconds-from-J2000 value.
\cpp{ConvertFrame(Real t\_tdb, ...)} is the important example
(\cref{chap:frame-conversions}).

That is an acceptable trade there rather than an oversight. A frame rotation is
a smooth function of time, and the Earth rotates at \SI{465}{\meter\per\second}
at the equator, so a \SI{245}{\nano\second} error in the time tag displaces a
position by ${\sim}\SI{0.1}{\milli\meter}$ --- four orders of magnitude below
the accuracy of the orientation model itself. The same
\SI{245}{\nano\second} interpreted as a \emph{light} time is \SI{74}{\meter},
which is not acceptable. The distinction is what the time is multiplied by.

\begin{lupntwarning}
Keep an instant in \cpp{Epoch}, and collapse to \code{Real} only at the
boundary of an interface that multiplies time by a velocity rather than
by $c$. Work that multiplies a time by $c$, or that differences two epochs,
must use \cpp{Epoch} or the offset accessors of
\cref{sec:tc-offset-accessors}.
\end{lupntwarning}

\subsection{Epoch series}

Much of the API is vectorized over epochs (\cpp{Propagate(x0, tfs)},
\cpp{GetBodyPosVel}, every \code{VEC\_DEF\_REAL} time function).
\cpp{EpochSeries} stores the split representation columnwise --- an exact
\code{int64} second per sample plus a \code{VecX} of fractions --- so a long
grid keeps the same guarantees a scalar \cpp{Epoch} does:

\begin{lstlisting}[language=cpp,caption={Building an exact uniform grid.},label={lst:tc-epochseries}]
// cpp/lupnt/conversions/epoch.h :: EpochSeries
static EpochSeries Linspace(const Epoch& start, Real step, int n);
VecX Since(const Epoch& ref) const;  // exact elapsed seconds relative to ref
VecX ToSeconds() const;              // LOSSY collapse, for interop
\end{lstlisting}

The $i$-th sample of \cpp{Linspace} is $\text{start} + i\,\text{step}$ computed
in the split representation, so a long grid does not accumulate the rounding
that repeated addition of a rounded step would. \cpp{Since(ref)} is the
precision-preserving way to feed a vectorized \code{Real} API: the large common
epoch is removed once, and the callee sees only small offsets.

\section{The conversion graph}
\label{sec:tc-graph}

A time scale is a coordinate on a particular four-dimensional reference
system, so converting between two of them means evaluating a relativistic
relationship rather than adding a constant \citep{turyshev2026cislunar}.
\lupnt{} represents those relationships as a directed graph whose edges return
the \emph{offset} between two scales,
%
\begin{equation}
\label{eq:tc-offset}
  \delta_{A\to B}(t) = t_B - t_A ,
\end{equation}
%
evaluated at the epoch's reading in scale $A$. A conversion finds the shortest
registered route (\cpp{numerics/graphs.h :: FindShortestPath}) and sums the
offsets along it,
%
\begin{equation}
\label{eq:tc-compose}
  \Delta_{A\to B}(t) = \sum_{k} \delta_{s_k \to s_{k+1}}\!\left(t_k\right),
  \qquad s_0 = A,\; s_n = B .
\end{equation}
%
Because \cref{eq:tc-compose} only ever sums small quantities, the total retains
full double precision, and it is the total --- not a converted absolute epoch
--- that is folded into $f$ by \cref{eq:tc-epochconv}.

\begin{figure}[htbp]
  \centering
  \begin{tikzpicture}[node distance=1.35cm and 1.15cm]
    \node[tsnode] (ut1) {UT1};
    \node[tsnode, right=of ut1] (utc) {UTC};
    \node[tsnode, right=of utc] (tai) {TAI};
    \node[tsnode, right=of tai] (tt)  {TT};
    \node[tsnode, right=of tt]  (tdb) {TDB};
    \node[tsnode, right=of tdb] (tcl) {TCL};
    \node[tsnode, right=of tcl] (lt)  {LT};
    \node[tsnode, below=of tai] (gps) {GPS};
    \node[tsnode, below=of tt]  (tcg) {TCG};
    \node[tsnode, below=of tdb] (tcb) {TCB};

    \draw[table]    (ut1) -- node[above, font=\tiny] {EOP}  (utc);
    \draw[table]    (utc) -- node[above, font=\tiny] {leap} (tai);
    \draw[constant] (tai) -- node[above, font=\tiny] {\SI{32.184}{\second}} (tt);
    \draw[model]    (tt)  -- node[above, font=\tiny] {model} (tdb);
    \draw[integral] (tdb) -- node[above, font=\tiny] {$\int$} (tcl);
    \draw[linear]   (tcl) -- node[above, font=\tiny] {$L_L$} (lt);
    \draw[constant] (tai) -- node[right, font=\tiny] {\SI{19}{\second}} (gps);
    \draw[linear]   (tt)  -- node[right, font=\tiny] {$L_G$} (tcg);
    \draw[linear]   (tdb) -- node[right, font=\tiny] {$L_B$} (tcb);

    \begin{scope}[shift={(0,-2.55)}, every node/.style={font=\scriptsize, anchor=west}]
      \draw[constant] (0,0)   -- (0.7,0);   \node at (0.8,0)   {constant};
      \draw[table]    (2.4,0) -- (3.1,0);   \node at (3.2,0)   {table};
      \draw[linear]   (4.2,0) -- (4.9,0);   \node at (5.0,0)   {linear};
      \draw[model]    (6.2,0) -- (6.9,0);   \node at (7.0,0)   {model};
      \draw[integral] (8.2,0) -- (8.9,0);   \node at (9.0,0)   {integral};
    \end{scope}
  \end{tikzpicture}
  \caption{The registered conversion graph. Every edge returns an
  \emph{offset}; a conversion is the sum of the offsets along the shortest
  route. Edge classes are ordered by cost: constant and table edges are
  arithmetic or a lookup, linear edges rescale elapsed time by
  $L\sim10^{-8}$, the TT\,$\leftrightarrow$\,TDB edge evaluates the active
  model, and the TDB\,$\leftrightarrow$\,TCL edge integrates from $T_0$. The
  graph is drawn undirected only for legibility: each edge is registered in
  both directions, with the inverse offset.}
  \label{fig:tc-graph}
\end{figure}

\begin{table}[htbp]
  \centering
  \caption{Edge classes and their cost. The spread is why the route matters.}
  \label{tab:tc-edges}
  \begin{tabularx}{\textwidth}{llL}
    \toprule
    Class & Edges & Cost \\
    \midrule
    constant & TAI\,$\leftrightarrow$\,TT, TAI\,$\leftrightarrow$\,GPS & arithmetic \\
    table    & TAI\,$\leftrightarrow$\,UTC (leap seconds),
               UTC\,$\leftrightarrow$\,UT1 (EOP) & table lookup \\
    linear   & TT\,$\leftrightarrow$\,TCG, TDB\,$\leftrightarrow$\,TCB,
               TCL\,$\leftrightarrow$\,LT & rescaling of elapsed time by $L\sim10^{-8}$ \\
    model    & TT\,$\leftrightarrow$\,TDB & Chebyshev fit, DE440 integral, or analytic series \\
    integral & TDB\,$\leftrightarrow$\,TCL & trapezoidal sweep from $T_0$ (1977) \\
    \bottomrule
  \end{tabularx}
\end{table}

TCL\,$\leftrightarrow$\,LT and TDB\,$\leftrightarrow$\,TCB are single linear
edges and are reached without touching the TDB\,$\leftrightarrow$\,TCL
integral. The search minimises hops, which coincides with minimum cost for
this edge set because the one expensive edge is also the only bridge to the
lunicentric scales: no route can avoid it that needs it, and no route incurs it
that does not. Adding an edge that breaks that coincidence would require a
cost-weighted search.

\begin{lstlisting}[language=cpp,caption={The registered edge table, abridged.},label={lst:tc-edges}]
// cpp/lupnt/conversions/epoch.cc :: TimeEdges
const std::map<std::pair<Time, Time>, TimeOffsetFn>& TimeEdges() {
  static const std::map<std::pair<Time, Time>, TimeOffsetFn> edges = {
      // --- constant offsets ---
      {{Time::TAI, Time::TT},  [](Real) { return Real(TT_TAI_OFFSET); }},
      {{Time::TAI, Time::GPS}, [](Real) { return Real(-19.0); }},

      // --- leap seconds / Earth orientation ---
      {{Time::TAI, Time::UTC},
       [](Real t_tai) { return Real(-GetTaiUtcDifference(TimeToMjd(t_tai).val())); }},
      {{Time::UTC, Time::UT1},
       [](Real t_utc) { return GetUt1UtcDifference(TimeToMjd(t_utc)); }},

      // --- linear rescalings (elapsed time only ever scaled by L ~ 1e-8) ---
      {{Time::TT,  Time::TCG},
       [](Real t_tt)  { return L_G / (1.0 - L_G) * (t_tt - T0Seconds()); }},
      {{Time::TDB, Time::TCB},
       [](Real t_tdb) { return (L_B * (t_tdb - T0Seconds()) - TDB_0) / (1.0 - L_B); }},
      {{Time::TCL, Time::LT},
       [](Real t_tcl) { return -L_L * (t_tcl - T0Seconds()); }},

      // --- the two model edges ---
      {{Time::TDB, Time::TT},  [](Real t_tdb) { return TtMinusTdb(t_tdb); }},
      {{Time::TDB, Time::TCL}, [](Real t_tdb) { return -TdbMinusTcl(t_tdb); }},
      // ... every edge is also registered in the reverse direction ...
  };
  return edges;
}
\end{lstlisting}

\subsection{Offset accessors}
\label{sec:tc-offset-accessors}

An interface that must \emph{return} an absolute epoch cannot benefit from
\cref{eq:tc-compose}; it is capped near the ULP however good the underlying
model is. \Cref{tab:tc-entrypoints} states the resulting contract.

\begin{table}[htbp]
  \centering
  \caption{What each entry point returns, and the accuracy that follows from
  the return type. Measured for TDB\,$\leftrightarrow$\,TT over 2020--2035.}
  \label{tab:tc-entrypoints}
  \begin{tabularx}{\textwidth}{lLl}
    \toprule
    Entry point & Returns & Accuracy \\
    \midrule
    \cpp{Epoch::To} / \cpp{TimeScaleOffset} & offset & \SI{4e-4}{\nano\second} \\
    \cpp{TtMinusTdb}, \cpp{TdbMinusTcl}, \cpp{TdbMinusLt} & offset & \SI{4e-4}{\nano\second} \\
    \cpp{ConvertTime} & absolute epoch & \SI{38}{\nano\second} rms, \SI{119}{\nano\second} peak \\
    \bottomrule
  \end{tabularx}
\end{table}

\cpp{ConvertTime(t, from, to)} delegates to \cpp{Epoch}; its floor is a
property of the return type, not of the model.

\begin{lstlisting}[language=py,caption={Choosing an interface in Python.},label={lst:tc-python}]
import pylupnt as pnt

e     = pnt.Epoch.from_seconds(t_tai, pnt.Time.TAI)
e_tdb = e.to(pnt.Time.TDB)                        # offset-composing, exact
dt    = pnt.time_scale_offset(e, pnt.Time.TCL)    # offset, sub-picosecond

t_tt  = pnt.convert_time(t_tai, pnt.Time.TAI, pnt.Time.TT)   # ~245 ns floor
\end{lstlisting}

\section{Proper versus coordinate time}
\label{sec:tc-proper-coordinate}

A clock measures \textbf{proper time} $\tau$ along its worldline, while a
\textbf{coordinate time} $t$ labels events in a four-dimensional relativistic
reference system (BCRS, GCRS, LCRS) and cannot be read off any single physical
clock \citep[Ch.~2.3.1]{iiyama2026thesis}. \cpp{IsCoordinateTimeScale}
partitions the enum: TT, TCG, TDB, TCB, TCL, and LT are coordinate scales;
UT1, UTC, TAI, and GPS are not.

\begin{lstlisting}[language=cpp]
// cpp/lupnt/conversions/time_conversions.cc :: IsCoordinateTimeScale
case Time::TT: case Time::TCG: case Time::TDB:
case Time::TCB: case Time::TCL: case Time::LT: return true;
\end{lstlisting}

\cpp{ConvertCoordinateTime} is the coordinate-only entry point. Its
position-aware overloads take a BCRS position $\vec{x}_{\mathrm{bcrs}}$ so that
the position-dependent terms of the TDB/TCB/TCL transforms are evaluated at the
event location rather than at the body centre.

\section{Atomic and terrestrial scales}
\label{sec:tc-atomic}

\paragraph{TAI to TT} is a defining constant offset
\citep[Eq.~2.38]{iiyama2026thesis}:
%
\begin{equation}
\label{eq:tc-tai-tt}
  \TT = \TAI + \SI{32.184}{\second} .
\end{equation}

\begin{lstlisting}[language=cpp]
// cpp/lupnt/conversions/time_conversions.cc :: TaiToTt / TtToTai
Real TaiToTt(Real t_tai) { return t_tai + TT_TAI_OFFSET; }   // TT_TAI_OFFSET = 32.184 s
Real TtToTai(Real t_tt)  { return t_tt  - TT_TAI_OFFSET; }
\end{lstlisting}

\paragraph{TAI to GPS} is the fixed \SI{19}{\second} offset that removes the
leap seconds accumulated at the GPS epoch
\citep[Eq.~2.39]{iiyama2026thesis}; GPS time runs without leap seconds
thereafter:
%
\begin{equation}
\label{eq:tc-gps}
  \GPST = \TAI - \SI{19}{\second} .
\end{equation}

\begin{lupntnote}
Only GPS is implemented among the GNSS scales. The thesis also lists
GLONASS $= \TAI - \SI{34}{\second}$ (leap-tracking, offset by UTC),
Galileo $= \TAI - \SI{19}{\second}$, and BeiDou $= \TAI - \SI{33}{\second}$
(Eqs.~2.40--2.42); these are not in the \code{Time} enum.
\end{lupntnote}

\paragraph{TAI and UTC} differ by the integer leap-second table
$\Delta_{\mathrm{AT}}(\mathrm{MJD})$ from \file{cpp/lupnt/interfaces/tai_utc.h}
\citep[Eq.~2.43]{iiyama2026thesis}:
%
\begin{equation}
\label{eq:tc-utc}
  \UTC = \TAI - \Delta_{\mathrm{AT}},
  \qquad
  \Delta_{\mathrm{AT}} = \cpp{GetTaiUtcDifference}(\mathrm{MJD}) .
\end{equation}

\paragraph{UTC and UT1} differ by the observed Earth-orientation correction
looked up from the EOP table (\file{cpp/lupnt/interfaces/eop.h}):
%
\begin{equation}
\label{eq:tc-ut1}
  \UTone = \UTC + (\UTone - \UTC),
  \qquad
  (\UTone - \UTC) = \cpp{GetUt1UtcDifference}(\mathrm{MJD}) .
\end{equation}

\begin{lstlisting}[language=cpp]
// cpp/lupnt/conversions/time_conversions.cc :: UtcToUt1
Real mjd_utc = t_utc / SECS_DAY + MJD_J2000_TT;
Real ut1_utc = GetUt1UtcDifference(mjd_utc);
return t_utc + ut1_utc;
\end{lstlisting}

UT1 in turn drives Earth rotation. \cpp{EarthRotationAngle} gives the ERA
%
\begin{equation}
\label{eq:tc-era}
  \theta = 2\pi\left(0.7790572732640
           + 1.00273781191135448\,\frac{t_{\UTone}}{86400}\right),
\end{equation}
%
and \cpp{GreenwichMeanSiderealTime} / \cpp{GreenwichApparentSiderealTime} give
GMST and GAST, the latter adding the equation of the equinoxes. These are the
only rotation quantities keyed on UT1; see \cref{chap:frame-conversions}.

\section{Barycentric and geocentric coordinate scales}
\label{sec:tc-bary}

The four Einstein-scaled coordinate times share the reference epoch
$T_0 = $ 1977-01-01 00:00:00 TAI (\code{MJD\_COORDINATE\_TT\_TCG\_TCB}) and the
defining rates \citep{iers2010}
%
\begin{equation}
\label{eq:tc-rates}
  L_B = 1.550519768\times10^{-8},\quad
  L_G = 6.969290134\times10^{-10},\quad
  \TDB_0 = \SI{-65.5}{\micro\second} .
\end{equation}
%
Writing $\tau = t - T_0$ for elapsed time from the coordinate-time origin, the
constant and linear edges of \cref{fig:tc-graph} are
%
\begin{align}
  \label{eq:tc-linear}
  \TCG - \TT   &= \frac{L_G}{1 - L_G}\,\tau_{\TT}, &
  \TCB - \TDB  &= \frac{L_B\,\tau_{\TDB} - \TDB_0}{1 - L_B}, \\
  \LT - \TCL   &= -L_L\,\tau_{\TCL}, &
  \UTC - \TAI  &= -\Delta_{\mathrm{AT}}(t) .
\end{align}

\begin{lstlisting}[language=cpp]
// cpp/lupnt/conversions/time_conversions.cc :: TtToTcg  (offset form)
Real tt_tcg = -L_G / (1.0 - L_G) * (mjd_tt - MJD_COORDINATE_TT_TCG_TCB) * SECS_DAY;
return t_tt - tt_tcg;

// cpp/lupnt/conversions/time_conversions.cc :: TcbToTdb
Real t_tdb = t_tcb - L_B * (mjd_tcb - MJD_COORDINATE_TT_TCG_TCB) * SECS_DAY + TDB_0;
\end{lstlisting}

Only the two remaining edges, TT\,$\leftrightarrow$\,TDB and
TDB\,$\leftrightarrow$\,TCL, require a relativistic model.

\section{The TT--TDB model}
\label{sec:tc-tttdb}

\subsection{The DE440 relation}

The barycentric edge is the DE440 relativistic relation
\citep[Eqs.~3--7]{park2021de440}. With $s_B = (1-L_G)/(1-L_B)$ and a station at
BCRS position $\vec{r}_S$,
%
\begin{equation}
\label{eq:tc-de440}
\begin{aligned}
  \TDB - \TT
    &= \frac{L_G - L_B}{1 - L_B}\,\tau_{\TDB}
     + s_B\,\TDB_0
     + \frac{s_B}{c^{2}} \int_{T_0 + \TDB_0}^{\TDB} f_2 \,\dd t
     - \frac{s_B}{c^{4}} \int_{T_0 + \TDB_0}^{\TDB} f_4 \,\dd t \\[2pt]
    &\quad + \frac{1}{c^{2}}\,\vec{v}_E \cdot (\vec{r}_S - \vec{r}_E)
     + \frac{1}{c^{4}}\left(3 w_{0E} + \tfrac{1}{2} v_E^{2}\right)
       \vec{v}_E \cdot (\vec{r}_S - \vec{r}_E),
\end{aligned}
\end{equation}
%
where the integrands collect the Earth's kinetic and potential terms,
%
\begin{align}
  \label{eq:tc-f2}
  f_2 &= \tfrac{1}{2} v_E^{2} + w_{0E} + w_{LE}, \\
  \label{eq:tc-f4}
  f_4 &= -\tfrac{1}{8} v_E^{4} - \tfrac{3}{2} v_E^{2} w_{0E}
         + 4\,\vec{v}_E \cdot \vec{w}_{iE}
         + \tfrac{1}{2} w_{0E}^{2} + \Delta_E ,
\end{align}
%
and $w_{0E}$ is the external potential at the Earth,
%
\begin{equation}
\label{eq:tc-w0e}
  w_{0E} = \sum_{i \neq E} \frac{GM_i}{r_{iE}}
  + U_{\mathrm{ring}}(r_{\odot E};\, GM_{\mathrm{ast}}, a_{\mathrm{ast}})
  + U_{\mathrm{ring}}(r_{\odot E};\, GM_{\mathrm{KB}}, a_{\mathrm{KB}}),
\end{equation}
%
with $w_{LE}$ the solar oblateness potential, $\vec{w}_{iE}$ the mass-weighted
velocity potential, and $\Delta_E$ the remaining $c^{-4}$ terms. The two ring
terms are \cref{sec:tc-smallbody}. At the geocentre $\vec{r}_S = \vec{r}_E$ and
both endpoint terms vanish, leaving only the integrals.

\subsection{The model chain}
\label{sec:tc-model-chain}

Evaluating \cref{eq:tc-de440} means sweeping from $T_0$, so it is not viable
per call. \cpp{TtMinusTdb(t)} is the single implementation of the TT--TDB
relation --- both \cpp{TtToTdb} and \cpp{TDBToTt} delegate to it, so the pair
round-trips exactly whichever model is active --- and it selects, in order:

\begin{enumerate}[leftmargin=*]
  \item \textbf{Chebyshev fit of the DE440t TT--TDB ephemeris.}
    \cpp{InitTtMinusTdbFit} samples the \file{de440t.bsp} time-ephemeris
    segment as a clean offset and fits piecewise Chebyshev segments (16-day
    segments, degree 12). Reproduces the kernel to \SI{4e-4}{\nano\second}.
  \item \textbf{DE440 integral}, \cref{eq:tc-de440}, opt-in via
    \cpp{SetTtTdbModel(TtTdbModel::DE440\_INTEGRAL)}. Trapezoidal sweep of the
    $c^{-2}$ and $c^{-4}$ integrands from $T_0$; agrees with the DE440t
    ephemeris to \SI{2.4}{\nano\second} rms with a residual secular term
    (\cref{sec:tc-boundaries}).
  \item \textbf{Auto-fit}: if no fit covers the requested epoch, one
    decade-wide window is built on demand (\cpp{SetTtTdbAutoFit}, default on).
    This is what makes the \emph{default} accuracy \SI{4e-4}{\nano\second} for
    every caller --- \cpp{ConvertTime}, \cpp{Epoch}, and the offset accessors
    alike.
  \item \textbf{Analytic fallback}, used only when the DE440t segment is
    unavailable: the two-term IAU/IERS series \citep{urban2013},
    %
    \begin{equation}
    \label{eq:tc-series}
      \TDB - \TT \approx
      \SI{0.001658}{\second}\,\sin g + \SI{0.000014}{\second}\,\sin 2g,
    \end{equation}
    %
    with $g = \ang{357.53} + \ang{0.9856003}\,d_{\mathrm{J2000}}$ the Earth's
    mean anomaly, good to $\lesssim\SI{2}{\milli\second}$.
\end{enumerate}

\begin{lupntnote}
The \textbf{position-aware} overloads \cpp{TtToTdb(t, x\_bcrs)} and
\cpp{TDBToTt(t, x\_bcrs)} evaluate the full GCRS\,$\leftrightarrow$\,BCRS
four-dimensional transform \citep[Eqs.~21--22]{turyshev2026cislunar}: a
trapezoidal integral of the $c^{-2}$ and $c^{-4}$ integrands built from the
Earth's barycentric velocity and the external Solar-System potential
$w_{\mathrm{ext}} = \sum_{B \neq E} GM_B / r_{EB}$, plus a position term
$-\vec{v}_E \cdot \vec{r}_E / c^2$. \cpp{TtToTdb(t, x\_bcrs)} inverts by
10-step Newton iteration.
\end{lupntnote}

\subsection{Small-body populations}
\label{sec:tc-smallbody}

DE440 integrates the solar system with 343 discrete asteroids, 30 individual
Kuiper-belt objects, and a 36-point circular ring representing the rest of the
Kuiper belt at \SI{44}{\au} \citep{park2021de440}. \lupnt{} carries no
ephemerides for any of them. Omitting that mass leaves a near-constant deficit
in the external potential, which \cref{eq:tc-de440} integrates into a purely
secular drift.

Each population is therefore added as a uniform circular ring of mass $M$ and
radius $a$, whose potential at heliocentric distance $r$ in the ring plane is
%
\begin{equation}
\label{eq:tc-ring}
  U(r) = \frac{2GM}{\pi(a+r)}\,K(k),
  \qquad k = \frac{2\sqrt{ar}}{a+r} .
\end{equation}
%
It is evaluated at the body's actual heliocentric distance at each step, so it
carries the small annual modulation rather than being held constant.

The factor $K$ is the complete elliptic integral of the first kind
\citep[Ch.~17]{abramowitz1964},
%
\begin{equation}
\label{eq:tc-elliptic}
  K(k) = \int_{0}^{\pi/2} \frac{\dd\theta}{\sqrt{1 - k^{2}\sin^{2}\theta}} .
\end{equation}
%
It arises for a geometric reason: the potential of a ring is the sum of $GM/d$
over every element of it, and the distance $d$ from the field point to a ring
element varies around the ring; integrating $1/d$ over the ring angle produces
exactly \cref{eq:tc-elliptic}. The modulus $k \in [0,1)$ measures how close the
field point is to the ring. As $k \to 0$ (far away) $K \to \pi/2$ and
\cref{eq:tc-ring} reduces to the point-mass potential $GM/r$; as $k \to 1$ (on
the ring) $K$ diverges logarithmically. For the Earth inside a \SI{44}{\au}
ring, $k$ is small and the ring is a mild correction to a point mass --- but a
correction with the right radial dependence, which is what the secular drift is
sensitive to.

\lupnt{} evaluates \cref{eq:tc-elliptic} by the arithmetic--geometric mean
(AGM) \citep[\S17.6]{abramowitz1964}: iterating
$a_{n+1} = \tfrac{1}{2}(a_n + b_n)$, $b_{n+1} = \sqrt{a_n b_n}$ from $a_0 = 1$,
$b_0 = \sqrt{1-k^{2}}$ converges quadratically to a common limit $M(k)$, and
$K(k) = \pi / \bigl(2 M(k)\bigr)$.

\begin{lstlisting}[language=cpp,caption={The ring potential, built from arithmetic and square roots so that it composes with automatic differentiation.},label={lst:tc-ring}]
// cpp/lupnt/conversions/time_conversions.cc :: RingPotential
// k' = sqrt(1-k^2). AGM converges quadratically (~5 iterations here) and is
// differentiable, so this composes with autodiff unlike std::comp_ellint_1.
Real RingPotential(Real r, double gm_ring, double a_ring) {
  if (gm_ring == 0.0) return Real(0.0);
  Real apr = a_ring + r;
  Real k2  = 4.0 * a_ring * r / (apr * apr);
  Real a = 1.0;
  Real b = sqrt(1.0 - k2);            // complementary modulus k'
  for (int i = 0; i < 12; i++) {
    Real a_next = 0.5 * (a + b);
    b = sqrt(a * b);
    a = a_next;
    if (abs((a - b).val()) < 1e-16) break;
  }
  Real K = PI / (2.0 * a);
  return 2.0 * gm_ring * K / (PI * apr);
}
\end{lstlisting}

\begin{table}[htbp]
  \centering
  \caption{Ring parameters, taken from DE440's own integration header rather
  than from independently published values. See the caveat below.}
  \label{tab:tc-rings}
  \begin{tabular}{lrr}
    \toprule
    Population & $GM$ [\si{\meter\cubed\per\second\squared}] & Ring radius [\si{\au}] \\
    \midrule
    Kuiper ring $+$ 30 KBOs & $1.1232\times10^{13}$ & 44 \\
    343 asteroids           & $1.7053\times10^{11}$ & 2.7 \\
    \bottomrule
  \end{tabular}
\end{table}

\begin{lupntwarning}
DE440 \emph{fitted} its ring mass during the integration, so
\cref{tab:tc-rings} is not the published Kuiper-belt total --- the
independently derived value is smaller by a factor of $1.43$. Only the value
DE440 actually integrated reproduces its TT--TDB. These are therefore
reproduction constants, not physical mass estimates: they should not be reused
as measurements of either population, nor tuned against any residual. For
sub-nanosecond absolute agreement with JPL use the Chebyshev fit instead.
\end{lupntwarning}

\section{Lunicentric coordinate scales}
\label{sec:tc-lunar}

The Moon's analogues follow \citet{turyshev2026cislunar}, with LCRS as the
lunar four-dimensional system. Reaching lunar time uses the same construction
one level down, with the Moon in place of the Earth.

\subsection{TCB to TCL}

TCL is the coordinate time of the Lunicentric Celestial Reference System. The
barycentric-to-lunicentric transform \citep[Eqs.~23,~25]{turyshev2026cislunar},
\citep[Eqs.~2.44--2.45]{iiyama2026thesis} is structurally identical to the
Earth case with the Moon substituted:
%
\begin{equation}
\label{eq:tc-tcbtcl}
  \TCL = \TCB
    - \frac{1}{c^2}\!\int\!\Bigl(\tfrac{1}{2}v_M^2 + w_{\mathrm{ext}}(\vec{x}_M)\Bigr)\dd\TCB
    - \frac{1}{c^4}\!\int (\cdots)\,\dd\TCB
    - \frac{\vec{v}_M \cdot \vec{r}_M}{c^2} - \cdots
\end{equation}
%
where $\vec{r}_M = \vec{x}_{\mathrm{bcrs}} - \vec{x}_M$,
$\vec{v}_M = \dd\vec{x}_M/\dd\TCB$, and
$w_{\mathrm{ext}}(\vec{x}_M) = \sum_{B \neq M} GM_B / |\vec{x}_M - \vec{x}_B|$
sums the Sun, Mercury, Venus, Earth, Mars, and the outer planets
(\code{kTclExternalBodies}), plus the same two rings of
\cref{sec:tc-smallbody}. Passing no position places the clock at the lunar
centre ($\vec{v}_M \cdot \vec{r}_M = 0$). \cpp{TclToTcb} inverts by Newton
iteration.

Equivalently, in the offset form the graph edge uses,
%
\begin{equation}
\label{eq:tc-tcl}
  \TDB - \TCL
   = -L_B\,\tau_{\TCB} + \TDB_0
   + \frac{1}{c^{2}} \int_{T_0}^{\TCB} g_2 \,\dd t
   + \frac{1}{c^{4}} \int_{T_0}^{\TCB} g_4 \,\dd t,
\end{equation}
%
\begin{equation}
\label{eq:tc-g24}
  g_2 = \tfrac{1}{2} v_M^{2} + w_{0M},
  \qquad
  g_4 = \tfrac{1}{8} v_M^{4} + \tfrac{3}{2} v_M^{2} w_{0M} - \tfrac{1}{2} w_{0M}^{2},
\end{equation}
%
where $w_{0M}$ is \cref{eq:tc-w0e} evaluated at the Moon.

\begin{lstlisting}[language=cpp]
// cpp/lupnt/conversions/time_conversions.cc :: TcbToTcl  (Turyshev Eq. 23/25)
auto [i2, i4] = IntegrateTcbToTclTerms(t_tcb);
return t_tcb - i2/(C*C) - i4/(C*C*C*C) + TcbToTclPositionTerm(t_tcb, x_bcrs);
\end{lstlisting}

\subsection{TCL to LT}

LT is the lunar analogue of TT: a defining rate on the selenoid
\citep[Eqs.~2.48--2.49]{iiyama2026thesis} with
$L_L = 3.139054\times10^{-11}$,
%
\begin{equation}
\label{eq:tc-lt}
  \LT = \TCL - L_L\,(\TCL - T_{L0}),
  \qquad
  \frac{\dd \LT}{\dd \TCL} = 1 - L_L .
\end{equation}

\begin{lstlisting}[language=cpp]
// cpp/lupnt/conversions/time_conversions.cc :: TclToLt  (T_L0 assumed = T_0)
Real lt_tcl = -L_L * (mjd_tcl - MJD_COORDINATE_TT_TCG_TCB) * SECS_DAY;
return t_tcl + lt_tcl;
\end{lstlisting}

\begin{lupntwarning}
$L_L$ is not yet internationally standardized. The code uses the Turyshev
selenoid value $3.139054\times10^{-11}$ (potential
$\Phi_L = 2.82123744381\times10^{6}\ \si{\meter\squared\per\second\squared}$),
whereas \citet{kopeikin2024lunar} adopt $3.14027\times10^{-11}$
\citep[Ch.~2.3.5]{iiyama2026thesis}. LT is therefore provisional pending
standardization.
\end{lupntwarning}

\subsection{Lunar time as a function of TT: the direct path}
\label{sec:tc-lt-direct}

Rather than chaining TT\,$\to$\,TDB\,$\to$\,TCB\,$\to$\,TCL\,$\to$\,LT,
\lupnt{} also implements LT directly as an Earth--Moon \emph{difference}
\citep[Eq.~57]{turyshev2026cislunar}, the TDB-argument form of thesis
Eq.~2.50, in \cpp{TdbToLtMinusTt}:
%
\begin{equation}
\label{eq:tc-ltmtt}
\begin{aligned}
  \LT - \TT
    &= \frac{L_G - L_L}{1 - L_B}\left(\TDB - T_0 - \TDB_0\right) \\[2pt]
    &\quad - \frac{1}{c^{2}} \int_{T_0 + \TDB_0}^{\TDB}
       \left[ \tfrac{1}{2} v_{EM}^{2}
            + \frac{GM_E - 2\,GM_M}{r_{EM}}
            + W^{\odot}_{EM}
            + \tfrac{3}{2} L_G \frac{GM_S}{r_{SE}} \right] \dd t \\[2pt]
    &\quad + \frac{1}{c^{2}}
       \Bigl[ \vec{v}_E \cdot \vec{r}_{EM} \Bigr]_{\TDB}^{T_0 + \TDB_0}
     - \frac{3}{c^{4}} \int_{T_0 + \TDB_0}^{\TDB}
       \frac{GM_S}{r_{SE}}\, \vec{v}_E \cdot \vec{v}_{EM} \,\dd t,
\end{aligned}
\end{equation}
%
with the solar $\ell = 2$ tidal term
%
\begin{equation}
\label{eq:tc-tidal}
  W^{\odot}_{EM} = \frac{3}{2}\,\frac{GM_S}{r_{SE}^{5}}
    \left[ (\vec{r}_{ES} \cdot \vec{r}_{EM})^{2}
         - \tfrac{1}{3} r_{SE}^{2} r_{EM}^{2} \right].
\end{equation}

\begin{lupntnote}
The factor of two on $GM_M$ in \cref{eq:tc-ltmtt} is not a typo: one $-GM_M$
comes from TCL\,--\,TCG, and a second from the $L_L$ rescaling that defines LT.
Because \cref{eq:tc-ltmtt} contains only Earth--Moon differences, distant
bodies --- including the small-body rings of \cref{sec:tc-smallbody} ---
cancel identically and are deliberately omitted from it. Including them there
would add ${\sim}10^{-9}\ \si{\meter\squared\per\second\squared}$ of numerical
noise rather than physics.
\end{lupntnote}

The integral is a trapezoidal sweep at $\dd t = \SI{0.01}{\dayunit}$. For
repeated queries in a window, \cpp{InitLtMinusTtFit} fits piecewise Chebyshev
segments (degree 12, 1-day segments) over the window in a single forward sweep
and caches them; \cpp{TdbToLtMinusTt} then uses the fast path
(\cpp{ChebyshevFitModel::Eval}, \file{cpp/lupnt/numerics/cheby_fit.h}).

\begin{lstlisting}[language=cpp]
// cpp/lupnt/conversions/time_conversions.cc :: TdbToLt / TtToLt
Real TdbToLt(Real t_tdb) { return TDBToTt(t_tdb) + TdbToLtMinusTt(t_tdb); }
Real TtToLt(Real t_tt)   { return TdbToLt(TtToTdb(t_tt)); }
\end{lstlisting}

\cpp{LtToTdb} and \cpp{LtToTt} invert by 10-step Newton iteration. That
\lupnt{} reaches LT by two genuinely different formulations --- the absolute
lunar potential of \cref{eq:tc-tcl} and the Earth--Moon difference of
\cref{eq:tc-ltmtt} --- makes their agreement a real check rather than a
tautology.

A separate lunar-surface proper-time term is provided by
\cpp{GetProperTimeCorrectionTcl(t\_tcg, x\_mci)}, the first-order clock
correction $-\vec{v}_{LE}\cdot(\vec{r}_x - \vec{r}_{LE})/c^2$ for a clock at
MCI position $\vec{x}_{\mathrm{mci}}$ relative to the Moon's geocentric motion.

\section{Exact calendar input}
\label{sec:tc-calendar}

The chain from a calendar date to an epoch is a second place where precision
can be lost before any conversion happens. Routing through a Modified Julian
Date rounds at the coarsest point in the chain: MJD counts days from 1858, so
one ULP of an MJD at present-day epochs is \SI{1129}{\nano\second} --- eight
times coarser than seconds-from-J2000, because its origin sits eight times
further away. Whole minutes survive that (their fraction-of-day is a dyadic
rational) but arbitrary clock times do not.

\cpp{Epoch::FromGregorian} therefore computes the integer second with integer
arithmetic \citep{hinnant2013} and carries only the sub-second remainder as a
floating-point value:

\begin{lstlisting}[language=cpp,caption={Integer-exact calendar conversion.},label={lst:tc-fromgregorian}]
// cpp/lupnt/conversions/epoch.cc :: Epoch::FromGregorian
const int64_t days = DaysFromCivil(year, month, day) - kDaysCivilToJ2000;
const double  s       = sec.val();
const double  s_floor = std::floor(s);
const int64_t whole   = days * 86400 + static_cast<int64_t>(hour) * 3600
                        + static_cast<int64_t>(min) * 60
                        + static_cast<int64_t>(s_floor) - kJ2000NoonOffset;
return Epoch(whole, sec - s_floor, scale);
\end{lstlisting}

The measured input error across a set of calendar dates including non-dyadic
clock times is \SI{0.0}{\nano\second}, against \SI{358}{\nano\second} for the
MJD-routed path. Dates before the Gregorian reform
defer to the MJD path, which handles the calendar discontinuity.

\section{Conversion of physical quantities and constants}
\label{sec:tc-scaling}

Because the scaled coordinate times (TT, TDB, LT) run at rates $(1-L_G)$,
$(1-L_B)$, $(1-L_L)$ relative to their unscaled partners (TCG, TCB, TCL),
spatial coordinates and gravitational parameters must be rescaled to keep $c$
and the equations of motion invariant
\citep[Eqs.~2.52--2.57]{iiyama2026thesis}:
%
\begin{align}
  \label{eq:tc-scale-tt}
  \mathcal{X}_{\TT}  &= (1-L_G)\,\mathcal{X}_{\TCG}, &
  \mu_{\TT}          &= (1-L_G)\,\mu_{\TCG}, \\
  \label{eq:tc-scale-tdb}
  \mathcal{X}_{\TDB} &= (1-L_B)\,\mathcal{X}_{\TCB}, &
  \mathcal{X}_{\LT}  &= (1-L_L)\,\mathcal{X}_{\TCL}.
\end{align}
%
\lupnt{} encodes these as the \code{CoordinateScale} factors $1 - L_x$, with
\cpp{CoordinateScaleRatio(from, to)} returning
$(1-L_{\mathrm{to}})/(1-L_{\mathrm{from}})$ for the rescale, guarded by
\cpp{AreCoordinateScalesConvertible} --- only
barycentric\,$\leftrightarrow$\,barycentric,
geocentric\,$\leftrightarrow$\,geocentric, and
lunicentric\,$\leftrightarrow$\,lunicentric pairs share a constant factor.

\begin{lstlisting}[language=cpp]
// cpp/lupnt/core/constants.h :: CoordinateScaleFactor
case CoordinateScale::TDB: return 1.0 - L_B;
case CoordinateScale::TT:  return 1.0 - L_G;
case CoordinateScale::TL:  return 1.0 - L_L;
\end{lstlisting}

\section{Model boundaries}
\label{sec:tc-boundaries}

\begin{itemize}[leftmargin=*]
  \item All API times are seconds from the per-scale J2000 origin; calendar
    input and output go through \cpp{GregorianToTime} /
    \cpp{TimeToGregorianString}, or --- preferably ---
    \cpp{Epoch::FromGregorian} (\cref{sec:tc-calendar}).
  \item Only GPS is implemented among the GNSS scales; GLONASS, Galileo, and
    BeiDou offsets are documented but not exposed.
  \item \cpp{ConvertTime} and \cpp{ConvertCoordinateTime} return absolute
    epochs and are therefore bounded near \SI{245}{\nano\second} regardless of
    model quality. Use \cpp{Epoch} or the offset accessors below that level.
  \item The DE440 integral of \cref{eq:tc-de440} carries a small secular
    difference against the DE440t ephemeris, dominated by the ring
    approximation of \cref{sec:tc-smallbody}. Refining the trapezoidal step
    from \SI{864}{\second} to \SI{216}{\second} moves the drift by
    \SI{0.0005}{\nano\second\per\yr}, so the residual is a model difference
    and not quadrature. Use the Chebyshev fit when absolute agreement with JPL
    matters.
  \item \cpp{TdbMinusTcl} and \cpp{TdbToLtMinusTt} integrate from $T_0$ per
    call. Both auto-fit a decade-wide Chebyshev window on demand
    (\cpp{SetTdbTclAutoFit}, \cpp{InitLtMinusTtFit}); with auto-fitting
    disabled a single lunar conversion costs seconds.
  \item $L_L$, and hence LT, is provisional pending international
    standardization.
\end{itemize}
